\medskip

Dans un repère orthonormé \Oij{} du plan, on considère les points A$(2~;~-1)$,
B$(0~;~3)$ et C$(3~;~1)$.
\begin{enumerate}
\item \begin{enumerate}
		\item Vérifier que $\vv{\text{AB}}\cdot\vv{\text{AC}} = 6$.


		\item Calculer $\left\|\vv{\text{AB}}\right\|$ et $\left\|\vv{\text{AC}}\right\|$, on donnera les valeurs exactes.



		\item Vérifier que $\cos(\widehat{\text{BAC}}) = 0,6$ et en déduire la mesure de l'angle $\widehat{\text{BAC}}$ au degré près.







		\end{enumerate}
\item 
	\begin{enumerate}
		\item Vérifier qu'une équation cartésienne de la droite (AB) est $2x + y - 3 = 0$.

















		\item On note H le pied la hauteur du triangle ABC issue du sommet C.

Déterminer les coordonnées du point H.






















	\end{enumerate}
\end{enumerate}

\vspace{0,5cm}
